\section{Bisimulation for Weaved models}

% \begin{definition}[\IML-bisimulation]
% 	Let $\wmodel_1 = \tuple{\Kwm[1], \Swm[1]}$ and $\wmodel_2 = \tuple{\Kwm[2], \Swm[2]}$ be two weaved models. Assume that for any $\kmodel_1,\kmodel_2 \in \Kwm[i]$ we have $Q_{\kmodel_1}$ disjoint $Q_{\kmodel_2}$. Define $\bigq[i] = \bigcup_{\kmodel \in \Kwm[i]}Q_{\kmodel}$, and then for $s \in \bigq[i]$ let 
% 	$\kmodel_s$ be the unique model in $\Kwm[i]$ that has $s$ in its domain. An \emph{\IML-bisimulation} between $\wmodel_1$ and $\wmodel_2$ is a non-empty relation $Z \subseteq \bigq[1] \times \bigq[2]$ satisfying the following properties.  Whenever $s Z s'$ then
% 	\begin{description}
%         \setlength{\itemsep}{0pt}
%         \item[\textnormal{(harmony)}] $V_{\kmodel_s}(s) = V_{\kmodel_{s'}}(s')$.
% 		\item[\textnormal{(zig)}] If $s R_{\kmodel_s} t$ then there is $t'$ such that $s' R_{\kmodel_{s'}} t'$ and $t Z t'$.
% 		\item[\textnormal{(zag)}] If $s' R_{\kmodel_{s'}} t'$ then there is $t$ such that $s R_{\kmodel_s} t$ and $t Z t'$.
% 		\item[\textnormal{(s-zig)}] If for some $\Sim$, $s \Sim t$ then there is $t'$ such that 
% 		from some $\Tim$, $s' \Tim t'$ and $t Z t'$.
% 		\item[\textnormal{(s-zag)}] If for some $\Sim$, $s' \Sim t'$ then there is $t$ such that 
% 		from some $\Tim$, $s \Tim t$ and $t Z t'$.
% 	\end{description}
% 	We write $\wmodel_1, s \bisim \wmodel_2, s'$ if there is a bisimulation $Z$ between 
% 	$\wmodel_1$ and $\wmodel_2$ such that $sZs'$.
% \end{definition}

\begin{definition}\label[definition]{def:wbisim}
	A \defn{bisimulation} between weaved models $\wmodel_1$ and $\wmodel_2$, written $Z \colon \wmodel_1 \bisim \wmodel_2$, is a
		non-empty relation $Z \subseteq \Qwm[1] \times \Qwm[2]$ such that,
			for all ${\pair {\kmodel_1} s} \in \Qwm[1]$ and all ${\pair {\kmodel_2} t} \in \Qwm[2]$
				with ${\pair {\kmodel_1} s} Z {\pair {\kmodel_2} t}$:
	\begin{description}
        \setlength{\itemsep}{0pt}
        \item[\textnormal{(harmony)}] $\Vkm[1](s) = \Vkm[2](t)$.
		\item[\textnormal{($\R$-zig)}]
			If there is $u \in \Wkm[1]$ with $s \Rkm[1] u$, then,
				there is $v \in \Wkm[2]$ such that $t \Rkm[2] u$ and ${\pair {\kmodel_1} u} Z {\pair {\kmodel_2} v}$.
		\item[\textnormal{($\R$-zag)}]
			If there is $v \in \Wkm[2]$ with $t \Rkm[2] v$, then,
				there is $u \in \Wkm[1]$ such that $s \Rkm[1] u$ and ${\pair {\kmodel_1} u} Z {\pair {\kmodel_2} v}$.
		\item[\textnormal{($\Fsymb$-zig)}]
			If there are ${\pair {\kmodel_3} u} \in \Qwm[1]$ and ${\Sim\colon \kmodel_1 \fsim \kmodel_3} \in \Swm[1]$ with $s \Sim u$, then,
				there are ${\pair {\kmodel_4} v} \in \Qwm[2]$ and ${\Tim\colon \kmodel_2 \fsim \kmodel_4} \in \Swm[2]$ with $t \Tim v$ and ${\pair {\kmodel_3} u} Z {\pair {\kmodel_4} v}$.
		\item[\textnormal{($\Fsymb$-zag)}]
			If there are ${\pair {\kmodel_4} v} \in \Qwm[2]$ and ${\Tim\colon \kmodel_2 \fsim \kmodel_4} \in \Swm[2]$ with $t \Tim v$, then,
				there are ${\pair {\kmodel_2} u} \in \Qwm[1]$ and ${\Sim\colon \kmodel_1 \fsim \kmodel_3} \in \Swm[1]$ with $s \Sim u$ and ${\pair {\kmodel_3} u} Z {\pair {\kmodel_4} v}$.
	\end{description}
	% We write $Z \colon {\pair {\wmodel_1} {\pair {\kmodel_1} s}} \bisim {\pair {\wmodel_2} {\pair {\kmodel_2} t}}$ if $Z$ is a bisimulation between $\wmodel_1$ and $\wmodel_2$ and ${\pair {\kmodel_1} s} Z {\pair {\kmodel_2} t}$.
\end{definition}

\begin{theorem}[Invariance]
	If $Z \colon \wmodel_1 \bisim \wmodel_2$ and ${\pair {\kmodel_1} s} Z {\pair {\kmodel_2} t}$, then, for all formulas $A$, we have $\xforces {\wmodel_1} {\pair {\kmodel_1} s} A$ iff $\xforces {\wmodel_2} {\pair {\kmodel_1} s} A$.
\end{theorem}
\begin{proof}
    \setcounter{ieqnnum}{0}
    The proof is by induction on $A$.

    \medskip
    \noindent
    The base case is: $\xforces {\wmodel_1} {\pair {\kmodel_1} s} p$ iff $\xforces {\wmodel_2} {\pair {\kmodel_2} t} p$.
	This follows immediately from the (harmony) condition, which guarantees that $\Vkm[1](s) = \Vkm[2](t)$, meaning $s$ and $t$ assign the same truth values to all proposition symbols.

    \medskip
    \noindent
    The inductive step proceeds with a case analysis.
    We prove some representative cases of the kind of reasoning involved.

	\medskip
    \noindent
    For diamonds, we prove $\xforces {\wmodel_1} {\pair {\kmodel_1} s} {\Diamond B}$ implies $\xforces {\wmodel_2} {\pair {\kmodel_2} t} {\Diamond B}$.
    Suppose $\xforces {\wmodel_1} {\pair {\kmodel_1} s} {\Diamond B}$.
    There is $u \in \Wkm[1]$ with \ieqn{s \Rkm[1] u} such that \ieqn{\xforces {\wmodel_1} {\pair {\kmodel_1} u} B}.
	Given \ieqnref{-1}, using the ($\R$-zig) condition, we know that there is \ieqn{v \in \Wkm[2]} with \ieqn{t \Rkm[2] v} such that ${\pair {\kmodel_1} u} Z {\pair {\kmodel_2} v}$.
	By the induction hypothesis, \ieqn{\xforces {\wmodel_2} {\pair {\kmodel_2} v} B}.
	Taken together, \ieqnref{-2}--\ieqnref0 imply $\xforces {\wmodel_2} {\pair {\kmodel_2} t} {\Diamond B}$ as required.
	The converse direction, $\xforces{\wmodel_2}{\pair{\kmodel_2}{t}}{\Diamond B}$ implies $\xforces{\wmodel_1}{\pair{\kmodel_1}{s}}{\Diamond B}$, is established by an analogous argument, this time applying the ($\R$-zag) condition in place of ($\R$-zig).

    \medskip
    \noindent
    For boxes, instead of proving $\xforces {\wmodel_1} {\pair {\kmodel_1} s} {\Box B}$ implies $\xforces {\wmodel_2} {\pair {\kmodel_2} t} {\Box B}$, we prove the contrapositive $\nxforces {\wmodel_2} {\pair {\kmodel_2} t} {\Box B}$ implies $\nxforces {\wmodel_1} {\pair {\kmodel_1} s} {\Box B}$.
    Suppose $\nxforces {\wmodel_2} {\pair {\kmodel_2} t} {\Box B}$.
    There are
		${\pair {\kmodel_4} u} \in \Qwm[2]$
			with \ieqn{{\pair {\kmodel_2} t} \Fwm[2] {\pair {\kmodel_4} u}}, and
		$v \in \Wkm[4]$
			with \ieqn{u \Rkm[4] v} 
	such that
		\ieqn{\nxforces {\wmodel_2} {\pair {\kmodel_4} v} B}.
	From \ieqnref{-2}, we know there is $\Tim\colon\kmodel_2 \fsim \kmodel_4 \in \wmodel_2$ such that $t \Tim u$.
	Given ${\pair {\kmodel_1} s} Z {\pair {\kmodel_2} t}$, the ($\Fsymb$-zag) condition ensures there are ${\pair {\kmodel_3} w} \in \Qwm[1]$ and $\Sim\colon\kmodel_1 \fsim \kmodel_3 \in \wmodel_1$ such that \ieqn{s \Sim w} and \ieqn{{\pair {\kmodel_3} w} Z {\pair {\kmodel_4} u}}.
	Having established \ieqnref{-3} and \ieqnref0, the ($\R$-zag) condition ensures there is $x \in \Wkm[3]$ with \ieqn{w \Rkm[3] x} such that \ieqn{{\pair {\kmodel_3} x} Z {\pair {\kmodel_4} v}}.
	By the inductive hypothesis, \ieqn{\nxforces {\wmodel_1} {\pair {\kmodel_3} x} {B}}.
	Taken together, \ieqnref{-4}, \ieqnref{-2}, and \ieqnref0, imply $\nxforces {\wmodel_1} {\pair {\kmodel_1} s} {\Box B}$ as required.
    The converse direction, $\xforces{\wmodel_2}{\pair{\kmodel_2}{t}}{\Box B}$ implies $\xforces{\wmodel_1}{\pair{\kmodel_1}{s}}{\Box B}$, is established by an analogous argument, this time applying the ($\Fsymb$-zig) and ($\R$-zig) conditions in place of ($\Fsymb$-zag) and ($\R$-zag), respectively.
\end{proof}

%%%% GENERATED WEAVED model

\begin{definition}\label[definition]{def:gen:wmodel}
    Let $\wmodel$ be a weaved model.
	For any $\kmodel \in \Kwm$, define the \defn{restriction} of $\wmodel$ to $\kmodel$ as a pair $\rwmodel = \tuple{\Krwm, \Srwm}$ where:
    \begin{enumerate}[leftmargin=\parindent]
        \setlength{\itemsep}{0pt}
        \item $\Krwm = \setof{\kmodel'}{\mbox{there is } {\Sim \colon \kmodel \fsim \kmodel'} \in \Swm}$; and
        \item $\Srwm = \Swm \cap (\Krwm \times \Krwm)$
    \end{enumerate}
    We also call $\rwmodel$ the \defn{$\kmodel$-generated weaved submodel} of $\wmodel$, or simply a generated weaved submodel if $\wmodel$ and $\kmodel$ are understood from the context.
\end{definition}

\begin{proposition}\label[proposition]{prop:bisim:rwmodel}
	Let $\wmodel$ be a weaved model.
	For all $\kmodel \in \Kwm$, the relation $Z = \setof{\pair {\pair {\kmodel'} s} {\pair {\kmodel'} s}}{ {\pair {\kmodel'} s} \in \Qof{\rwmodel}}$ is a bisimulation between $\rwmodel$ and $\wmodel$.
\end{proposition}

\begin{corollary}\label[corollary]{prop:rwmodel}
    For all formulas $A$, we have $\wforces{\pair \kmodel s}{A} \mbox{ iff } \xforces{\rwmodel}{\pair \kmodel s}{A}$.
\end{corollary}

%%%% UNRAVELING WEAVED model

\begin{definition}\label[definition]{def:unraveling}
	Let $\wmodel$ be a weaved model and let $\kmodel \in \Kwm$.
	To properly define the unraveling of $\wmodel$ around $\kmodel$, we first need to introduce a few preliminary notions.
        \lutz{that's not a definition} \lutz{and never write a definition that spans more than a page}

	\medskip
	\noindent
	By a \defn{simulation thread} starting at $\kmodel$ we mean a finite sequence $\pi = \Sim_1 \dots \Sim_n$ of simulations in $\Swm$ such that: $\dom(\Sim_1) = \kmodel$, and the composition ${\Sim_1 \dots \Sim_n \in \Swm}$.
	We take the empty sequence, denoted by $\varepsilon$, to stand for the starting model $\kmodel$.
	We use $|\pi|$ to indicate the length of a thread $\pi$.
	For each $i \in [1{:}|\pi|]$, we let $\pi_i$ be the codomain of the $i$-th simulation in $\pi$, that is, $\pi_i = \cod(\Sim_i) = \kmodel_i$, with the convention $\pi_0 = \kmodel$.
	We write $\pi[{:}i]$ to denote the initial segment of $\pi$ up to the $i$-th simulation, that is, $\pi[{:}i] = \Sim_1 \dots \Sim_i$, with the convention $\pi[{:}0] = \varepsilon = \kmodel$.
	Finally, we use $\Pi$ to denote the collection of all such simulation threads in $\wmodel$.
	Intuitively, simulation threads make paths in the weaved model explicit: they record step by step how one moves through the model. As we will show below, simulation threads provide the fundamental building blocks for the unraveling construction.
	
	\medskip
	\noindent
	To transform abstract simulation threads into concrete components of the unraveling, we tie each thread to a unique copy of the Kripke model it reaches. To achieve this, we associate to a each $\pi \in \Pi$ a set $\Qof{\pi} = \setof{s_{\pi}}{s \in \Qof{\pi_{|\pi|}}}$ and a bijection $f_{\pi}\colon {\Qof{\pi_{|\pi|}} \bsym \Qof{\pi}}$ such that $f_{\pi}(s) = s_\pi$.
	This definition of $f_{\pi}$ provides us with a systematic way of generating the unique Kripke models that will play the role of the nodes of tree-like model characteristic of unravelings.

	\medskip
	\noindent
	Concretely, every $f_{\pi}$ yields a Kripke model $\kmodel_{\pi} = \tuple{\Qof{\pi}, \Rof{\pi}, \Vof{\pi}}$ where:
	\begin{enumerate}[leftmargin=\parindent,resume]
        \setlength{\itemsep}{0pt}
		% \item $\Qof{\pi} = \cod(f_{\pi})$;
		\item $\Rof{\pi} = \setof{\pair {s_\pi} {t_\pi}}{ s \Rof{\pi_{|\pi|}} t}$; and
		\item for all $s_{\pi} \in \Qof{\pi}$, $\Vof{\pi}(s_{\pi}) = \Vof{\pi_{|\pi|}}(s)$.
	\end{enumerate}
	It is immediate from the definition of $f_{\pi}$ that each $\kmodel_{\pi}$ is indeed a unique isomorphic copy of $\pi_{|\pi|}$.

	\medskip
	\noindent
	Having identified the Kripke models that will appear in the unraveling, we turn to defining the simulations that will connect these models together. Intuitively, these simulations must branch in a disciplined manner so that, taken as a whole, they generate the tree-like shape characteristic of unravelings.
	The key lies in the notion of an extension of a simulation thread.
	
	\medskip
	\noindent
	We say that a simulation thread $\pi'$ is an \defn{extension} of a simulation thread $\pi$ iff $\pi = \pi'[{:}|\pi|]$. We write $\pi\pi'$ if $\pi'$ is an extension of $\pi$. Extensions provide the natural mechanism for controlling how simulations branch in the unraveling.

	\medskip
	\noindent
	Concretely, for every pair $\pi\pi'$ with $\pi' = \Sim_1 \dots \Sim_{|\pi|} \Sim_{(|\pi|+1)} \dots \Sim_{|\pi'|}$, we define
	\begin{enumerate}[leftmargin=\parindent,resume]
        \setlength{\itemsep}{0pt}
		\item $\Sim_{\pi\pi'} = f_{\pi}^{-1} \Sim_{(|\pi|+1)} \dots \Sim_{|\pi'|} f_{\pi'}$.
	\end{enumerate}
	It follows directly from the definition that $\Sim_{\pi\pi'}$ is a simulation from $\kmodel_{\pi}$ to $\kmodel_{\pi'}$.
	Composition of these simulations is defined in the natural way: for any $\pi, \pi', \pi'' \in \Pi$, we set $\Sim_{\pi\pi'}\Sim_{\pi'\pi''} = \Sim_{\pi\pi''}$. This composition is clearly well-defined.
	In addition, for every $\pi \in \Pi$, there exists a simulation $\Sim_{\pi\pi}$ that behaves as an identity: for all $\pi' \in \Pi$, if $\pi'$ is an extension of $\pi$, then, $\Sim_{\pi\pi}\Sim_{\pi\pi'} = \Sim_{\pi\pi'}$; and if $\pi$ is an extension of $\pi'$, then, $\Sim_{\pi'\pi}\Sim_{\pi\pi} = \Sim_{\pi'\pi}$.
	Thus, $\Sim_{\pi\pi}$ serves as the identity simulation for $\kmodel_{\pi}$.

	\medskip
	\noindent
	Having set up the necessary definitions, we can now put them together into the unraveling of $\wmodel$ around $\kmodel$. This construction turns the abstract machinery developed above into a tree-like model that makes explicit how simulation threads unfold.
	Precisely, we define the \defn{unraveling} of $\wmodel$ around $\kmodel$ as a pair $\uwmodel = \tuple{\Kuwm, \Suwm}$ where:
	\begin{enumerate}[leftmargin=\parindent,resume]
        \setlength{\itemsep}{0pt}
		\item $\Kuwm = \setof{\kmodel_\pi}{\pi \in \Pi}$;
		\item $\Suwm = \setof{\Sim_{\pi\pi'}}{\pi,\pi' \in \Pi \mbox{ and } \pi' \mbox{ is an extension of } \pi}$.
	\end{enumerate}
\end{definition}

\begin{proposition}
	The unraveling of a weaved model $\wmodel$ around a Kripke model $\kmodel \in \Kwm$ is a weaved model.
\end{proposition}
% \begin{proof}
% 	\setcounter{ieqnnum}{0}
% 	By construction, each $\kmodel_\pi \in \Kuwm$ is a Kripke model.
% 	The non-trivial part is to prove that each $\Sim^i_\pi \in \Suwm$ is a seesaw simulation from $\kmodel_{\pi[i]}$ to $\kmodel_\pi$. %, i.e., $\Sim^i_\pi \colon \kmodel_{\pi[i]} \fsim \kmodel_\pi$.
% 	To establish this, we must verify:
% 	\ieqn{\Sim^i_\pi \subseteq \Qof{\pi[i]} \times \Qof{\pi}}; and
% 	\ieqn{\Sim^i_\pi} satisfies conditions (add), (see), and (saw) of seesaw simulations.
% 	Condition \ieqnref{-1} follows directly from the definition of $\Sim^i_\pi$.
% 	We establish $\Sim^i_\pi$ is a simulation below.

% 	\medskip
% 	\noindent
% 	To simplify notation, let
% 		\ieqn{\pi = \kmodel_1 \Sim_1 \dots \Sim_{(n-1)} \kmodel_n}.
% 	Moreover, let
% 		\ieqn{\Sim = (\Sim_i \fcomp \dots \Sim_n)\colon \kmodel_i \fsim \kmodel_n \in \Swm}.

% 	\medskip
% 	\noindent
% 	For (add), choose an arbitrary $s \Sim^i_\pi t$.
% 	We must prove $\Vof{\pi[i]}(s) \subseteq \Vof{\pi}(t)$.
% 	By definition $\Vof{\pi[i]}(s) = \Vkm[i](s)$, and $\Vof{\pi}(t) = \Vkm[n](t)$.
% 	The result follows directly from the (add) condition in \ieqnref0.
	
% 	\medskip
% 	\noindent
% 	For (see), again choose an arbitrary $s \Sim^i_\pi t$.
% 	By definition, from \ieqnref0, we have $s_i \Sim t_n$.
% 	Take $u \in \Qof{\pi[i]}$ such that $s \Rof{\pi[i]} u$.
% 	This scenario is best illustrated by the following arrow-chasing diagram, which displays how the prefix $\pi[i]$ aligns with the full thread $\pi$ via the simulations~in~$\Sim$:
% 	\[
% 		\begin{tikzcd}
% 			s_1
% 				\arrow[r,dotted, "\Sim_1"]
% 				\arrow[d, swap, "{\Rkm[1]}" {yshift=1pt}]
% 			% &
% 			% s_2
% 			% 	\arrow[r,dotted, "\Sim_2" {xshift=-2pt}]
% 			% 	\arrow[d, "\Rof{\pi_2}" {yshift=1pt}]
% 			&
% 			\dots
% 				\arrow[r,dotted, "\Sim_{(i-1)}" {xshift=-2pt}]
% 			&
% 			s_i
% 				\arrow[r,dotted, "\Sim_i"]
% 				\arrow[rr, bend left=35,dotted, "\Sim"]
% 				\arrow[d, swap, "{\Rkm[i]}" {yshift=1pt}]
% 			% &
% 			% s_{(i+1)}
% 			% 	\arrow[r,dotted, "\Sim_{(i+1)}" {xshift=-2pt}]
% 			% 	\arrow[d, "\Rof{\pi_{(i+1)}}" {yshift=1pt}]
% 			&
% 			\dots
% 				\arrow[r,dotted, "\Sim_{(n-1)}" {xshift=-2pt}]
% 			&
% 			t_{n}
% 				\arrow[d, "{\Rkm[n]}" {yshift=1pt}]
% 			\\
% 			u_1
% 				\arrow[r, swap, dotted, "\Sim_1"]
% 			% &
% 			% u_2
% 			% 	\arrow[r, swap, dotted]{}{\Sim_2}
% 			&
% 			\dots
% 				\arrow[r, swap, dotted, "\Sim_{(i-1)}" {xshift=-2pt}]
% 			&
% 			u_i
% 				\arrow[r, swap, dotted, "\Sim_i"]
% 				\arrow[rr, bend left=35,dotted, "\Sim"]
% 			% &
% 			% u_{(i+1)}
% 			% 	\arrow[r, swap, dotted]{}{\Sim_{(i+1)}}
% 			&
% 			\dots
% 				\arrow[r, swap, dotted, "\Sim_{(n-1)}" {xshift=-2pt}]
% 			&
% 			u_{n}
% 		\end{tikzcd}
% 	\]
% 	Formally, $s \Rof{\pi[i]} u$ ensures that the tuple $s_1,\dots,s_i$ is linked step by step to $u_1,\dots,u_i$ through the chain of simulations.
% 	Then, $s \Sim^i_\pi t$ implies, using the (see) condition for $\Sim$, that the tuple $s_1,\dots,s_i,\dots,s_n$ extends coherently along the entire thread, yielding the sequence $v = u_1 \dots u_n \in \Qof{\pi}$.
% 	Consequently, we obtain both $t \Rof{\pi} v$ and $u \Sim^i_\pi v$, exactly as required for condition (see).
% 	The argument for (saw) is entirely analogous using the (saw) condition in $\Sim$ instead of (see) in $\Sim$.

% 	\medskip
% 	\noindent
% 	To conclude the proof we must prove that $\Suwm$ is closed under identities and composition.
% 	It is clear that $\Sim^1_\pi \in \Suwm$ acts as the identity on $\kmodel_\pi$.
% 	For composition, we note that $\Sim^i_{\pi}\colon \kmodel_{\pi[i]} \fsim \kmodel_{\pi}$ and $\Tim^j_{\pi'}\colon \kmodel_{\pi'[j]} \fsim \kmodel_{\pi'}$ can be composed only if $\pi = \pi'[j]$.
% 	Then, we define $s (\Sim^i_{\pi} \fcomp \Tim^j_{\pi'}) t$ iff the following conditions are met:
% 	\begin{enumerate}[leftmargin=\parindent,label=(\alph*)]
%         \setlength{\itemsep}{0pt}
% 		\item $s \in \Qof{\pi[i]}$;
% 		\item $t \in \Qof{\pi'}$; and
% 		\item exists $u \in \Qof{\pi}$ such that $s_i (\Sim_i \fcomp \dots \fcomp \Sim_{(|\pi|-1)}) u_{|\pi|} (\Tim_j \fcomp \dots \fcomp \Tim_{|\pi'|}) t_{|\pi'|}$.
% 	\end{enumerate}
% \end{proof}

\begin{proposition}\label[proposition]{prop:bisim:uwmodel}
	Let $\rwmodel$ be a $\kmodel$-generated weaved model.
	Then, for all
		$\pair {\kmodel'} s \in \Qof{\rwmodel}$ and all
		$\pair {\kmodel_\pi} t \in \Qof{\uwmodel}$,
	define
		\[
			{{\pair {\kmodel'} s} Z {\pair {\kmodel_{\pi}} {s_{\pi}}}}
			\quad
			\mbox{iff}
			\quad
			{\kmodel' = \pi_{|\pi|}}
		\]
	It follows that $Z$ is a bisimulation between $\rwmodel$ and $\uwmodel$.
\end{proposition}
\begin{proof}
	The key observation is that $Z$ provides a pointwise correspondence between the Kripke models in the generated weaved model and their isomorphic copies in the unraveling. This correspondence immediately guarantees the ($\R$-zig) and ($\R$-zag) conditions of a bisimulation.

	\medskip
	\noindent
	In turn, consider ${\Sim \colon \kmodel' \fsim \kmodel''} \in \Srwm$. Then, extending $\pi$ by $\Sim$ yields a simulation thread $\pi' = \pi\Sim \in \Pi$. This extension ensures the unraveling includes the simulation ${\Sim_{\pi\pi'} \colon \kmodel_\pi \fsim \kmodel_{\pi'}} \in \Suwm$, which guarantees ($\Fsymb$-zig).

	\medskip
	\noindent
	Conversely, given ${\Tim \colon \kmodel_{\pi'} \fsim \kmodel_{\pi''}} \in \Suwm$, we can move back to the weaved model: the bijections $f_{\pi'}$ and $f_{\pi''}$ ensure the existence of the simulation ${\Sim = f_{\pi'} \Tim f_{\pi''}^{-1} \colon \kmodel' \fsim \pi''_{|\pi''|}} \in \Srwm$, which provides ($\Fsymb$-zag).
\end{proof}

\begin{corollary}\label[corollary]{prop:uwmodel}
	For all formulas $A$, $\xforces{\rwmodel}{\pair \kmodel s}{A}$ iff $\xforces{\uwmodel}{\pair {\kmodel_{\varepsilon}} {s_{\varepsilon}}}{A}$.
\end{corollary}

\begin{definition}\label[definition]{def:idepth}
	The \defn{intuitionistic degree} of a formula $A$ is defined as:
	\begin{align*}
		\ideg(p) & = 0 \\
		\ideg(\bot) & = 0 \\
		\ideg(B \lor C) &= \max(\ideg(B),\ideg(C)) \\
		\ideg(B \land C) &= \max(\ideg(B),\ideg(C)) \\
		\ideg(B \to C) &= 1+\max(\ideg(B),\ideg(C)) \\
		\ideg(\Diamond B) &= \ideg(B) \\
		\ideg(\Box B) &= 1+ \ideg(B)
	\end{align*}
\end{definition}

\begin{definition}\label[definition]{def:nbisimulation}
	Let $\wmodel_1$ and $\wmodel_2$ be weaved models.
	Suppose we are given a finite sequence of non-empty relations $Z^n \subseteq Z^{(n-1)} \subseteq \dots \subseteq Z^0$ on $\Qwm[1] \times \Qwm[2]$, where each $Z^i$ satisfies the conditions (harmony), ($\R$-zig), and ($\R$-zag) from \Cref{def:wbisim}.
	We say that $Z^n$ is an \defn{$n$-bisimulation} between $\wmodel_1$ and $\wmodel_2$, and write $Z^n \colon \wmodel_1 \bisim^n \wmodel_2$, iff for all $i \in [1{:}n]$, and all ${\pair {\kmodel_1} s} \in \Qwm[1]$ and ${\pair {\kmodel_2} t} \in \Qwm[2]$ with ${\pair {\kmodel_1} s} Z^i {\pair {\kmodel_2} t}$, the following holds:
	\begin{description}
        \setlength{\itemsep}{0pt}
		\item[\textnormal{($\Fsymb^i$-zig)}]
			If there are ${\pair {\kmodel_3} u} \in \Qwm[1]$ and ${\Sim\colon \kmodel_1 \fsim \kmodel_3} \in \Swm[1]$ with $s \Sim u$, then,
				exist ${\pair {\kmodel_4} v} \in \Qwm[2]$ and ${\Tim\colon \kmodel_2 \fsim \kmodel_4} \in \Swm[2]$ with $t \Tim v$ and ${\pair {\kmodel_3} u} Z^{(i-1)} {\pair {\kmodel_4} v}$.
		\item[\textnormal{($\Fsymb^i$-zag)}]
			If there are ${\pair {\kmodel_4} v} \in \Qwm[2]$ and ${\Tim\colon \kmodel_2 \fsim \kmodel_4} \in \Swm[2]$ with $t \Tim v$, then,
				exist ${\pair {\kmodel_2} u} \in \Qwm[1]$ and ${\Sim\colon \kmodel_1 \fsim \kmodel_3} \in \Swm[1]$ with $s \Sim u$ and ${\pair {\kmodel_3} u} Z^{(i-1)} {\pair {\kmodel_4} v}$.
	\end{description}
\end{definition}

\begin{theorem}[$n$-invariance]\label[theorem]{prop:ninv}
	If $Z^n \colon \wmodel_1 \bisim^n \wmodel_2$ and ${\pair {\kmodel_1} s} Z^n {\pair {\kmodel_2} t}$, then, for all formulas $A$ with $\ideg(A) \leq n$, we have $\xforces {\wmodel_1} {\pair {\kmodel_1} s} A$ iff $\xforces {\wmodel_2} {\pair {\kmodel_1} s} A$.
\end{theorem}
\begin{proof}
    \setcounter{ieqnnum}{0}
    The proof is by induction on two measures: the degree and the size of the formula $A$.
	Precisely, we associate with each formula $A$ a complexity pair $\pair {\deg(A)} {\size(A)}$, the proof is by strong induction on the lexicographic order on such pairs.

    \medskip
    \noindent
    For the base case we need to establish $\xforces {\wmodel_1} {\pair {\kmodel_1} s} p$ iff $\xforces {\wmodel_2} {\pair {\kmodel_2} t} p$.
	This follows directly from the (harmony) condition in $Z^0$, which guarantees that $\Vkm[1](s) = \Vkm[2](t)$, meaning $s$ and $t$ assign the same truth values to all proposition symbols.

    \medskip
    \noindent
    For the inductive step, we choose a pair ${\pair {\kmodel_1} s} Z^{(i+1)} {\pair {\kmodel_2} t}$, and a formula $A$ with complexity $(i+1,j)$. The inductive hypothesis tells us that the result holds for all pairs ${\pair {\kmodel_3} u} Z^i {\pair {\kmodel_4} v}$ and all formulas with complexity $(k,\ell) < (i+1,j)$. We are required to prove that the result holds for $A$. We proceed with a case analysis.

	\medskip
    \noindent
    For conjunctions, we prove $\xforces{\wmodel_1}{\pair{\kmodel_1}{s}}{B \land C}$ implies $\xforces{\wmodel_2}{\pair{\kmodel_2}{t}}{B \land C}$.
	Suppose $\xforces{\wmodel_1}{\pair{\kmodel_1}{s}}{B \land C}$.
	By the semantics of conjunction, this means that $\xforces{\wmodel_1}{\pair{\kmodel_1}{s}}{B}$ and $\xforces{\wmodel_1}{\pair{\kmodel_1}{s}}{C}$.
	The complexity pairs associated with $B$ and $C$ both precede $(i+1,j)$ in the lexicographic order.
	Therefore, by the inductive hypothesis, we obtain $\xforces{\wmodel_2}{\pair{\kmodel_2}{t}}{B}$ and $\xforces{\wmodel_2}{\pair{\kmodel_2}{t}}{C}$. Together, these yield $\xforces{\wmodel_2}{\pair{\kmodel_2}{t}}{B \land C}$.
	The converse direction, $\xforces{\wmodel_2}{\pair{\kmodel_2}{t}}{B \land C}$ implies $\xforces{\wmodel_1}{\pair{\kmodel_1}{s}}{B \land C}$, follows by an entirely symmetric argument.
	The case for disjunctions proceeds in exactly the same way.

	\medskip
    \noindent
    For diamonds, we prove $\xforces {\wmodel_1} {\pair {\kmodel_1} s} {\Diamond B}$ implies $\xforces {\wmodel_2} {\pair {\kmodel_2} t} {\Diamond B}$.
    Suppose $\xforces {\wmodel_1} {\pair {\kmodel_1} s} {\Diamond B}$.
    There is $u \in \Wkm[1]$ with \ieqn{s \Rkm[1] u} such that \ieqn{\xforces {\wmodel_1} {\pair {\kmodel_1} u} B}.
	Given \ieqnref{-1}, using the ($\R$-zig) condition in $Z^{(i+1)}$, we know that there is \ieqn{v \in \Wkm[2]} with \ieqn{t \Rkm[2] v} such that ${\pair {\kmodel_1} u} Z^{(i+1)} {\pair {\kmodel_2} v}$.
	Now, the complexity pair associated with $B$ precedes $(i+1,j)$ in the lexicographic order.
	Therefore, by the induction hypothesis, \ieqn{\xforces {\wmodel_2} {\pair {\kmodel_2} v} B}.
	Taken together, \ieqnref{-2}--\ieqnref0 imply $\xforces {\wmodel_2} {\pair {\kmodel_2} t} {\Diamond B}$ as required.
	The converse direction, $\xforces{\wmodel_2}{\pair{\kmodel_2}{t}}{\Diamond B}$ implies $\xforces{\wmodel_1}{\pair{\kmodel_1}{s}}{\Diamond B}$, is established by an analogous argument, this time using ($\R$-zag) in place of ($\R$-zig).

    \medskip
    \noindent
    For boxes, instead of proving $\xforces {\wmodel_1} {\pair {\kmodel_1} s} {\Box B}$ implies $\xforces {\wmodel_2} {\pair {\kmodel_2} t} {\Box B}$, we prove the contrapositive $\nxforces {\wmodel_2} {\pair {\kmodel_2} t} {\Box B}$ implies $\nxforces {\wmodel_1} {\pair {\kmodel_1} s} {\Box B}$.
    Suppose $\nxforces {\wmodel_2} {\pair {\kmodel_2} t} {\Box B}$.
    There are
		${\pair {\kmodel_4} u} \in \Qwm[2]$ with
			\ieqn{{\pair {\kmodel_2} t} \Fwm[2] {\pair {\kmodel_4} u}}, and
		$v \in \Wkm[4]$ with
			\ieqn{u \Rkm[4] v} 
		such that
			\ieqn{\nxforces {\wmodel_2} {\pair {\kmodel_4} v} B}.
	From \ieqnref{-2}, exists $\Tim\colon\kmodel_2 \fsim \kmodel_4 \in \wmodel_2$ such that $t \Tim u$.
	Since ${\pair {\kmodel_1} s} Z^{(i+1)} {\pair {\kmodel_2} t}$, the ($\Fsymb^{(i+1)}$-zag) condition ensures there are
		${\pair {\kmodel_3} w} \in \Qwm[1]$ and
		$\Sim\colon\kmodel_1 \fsim \kmodel_3 \in \wmodel_1$
		such that
			\ieqn{s \Sim w} and
			\ieqn{{\pair {\kmodel_3} w} Z^i {\pair {\kmodel_4} u}}.
	Now, \ieqnref{-3} and \ieqnref0 imply, by the ($\R$-zag) condition in $Z^i$,
		there are
			$x \in \Wkm[3]$
			such that
				\ieqn{w \Rkm[3] x} and
				\ieqn{{\pair {\kmodel_3} x} Z^i {\pair {\kmodel_4} v}}.
	Since $\ideg(B) = i$, by the induction hypothesis, we obtain
		\ieqn{\nxforces{\wmodel_1}{\pair {\kmodel_3} {x}}{B}}.
	Taken together, \ieqnref{-4}, \ieqnref{-2}, and \ieqnref0, imply $\nxforces {\wmodel_1} {\pair {\kmodel_1} s} {\Box B}$ as required.
    The converse direction, $\xforces{\wmodel_2}{\pair{\kmodel_2}{t}}{\Box B}$ implies $\xforces{\wmodel_1}{\pair{\kmodel_1}{s}}{\Box B}$, is established by an analogous argument, this time applying the ($\Fsymb$-zig) and ($\R$-zig) conditions in place of ($\Fsymb$-zag) and ($\R$-zag), respectively.
\end{proof}

\begin{definition}\label[definition]{def:prunning}
	Let $\uwmodel$ be the unraveling of a weaved model $\wmodel$ around a Kripke model $\kmodel \in \Kwm$.
	The \defn{prunning} of $\uwmodel$ at height $n$ is a weaved model $\pwmodel = \tuple{\Kpwm, \Spwm}$ where:
	\begin{enumerate}[leftmargin=\parindent]
		\setlength{\itemsep}{0pt}
		\item For all $\kmodel_{\pi} \in \Kuwm$, $|\pi| \leq n$ implies $\kmodel_{\pi} \in \Kpwm$;
		\item For all $\Sim\colon\kmodel_{\pi}\fsim\kmodel_{\pi'} \in \Suwm$, if $\kmodel_{\pi},\kmodel_{\pi'} \in \Kpwm$, then, $\Sim \in \Sbwm$.
	\end{enumerate}
\end{definition}

\begin{proposition}\label[proposition]{prop:bounded:height}
	Let $\uwmodel$ be the unraveling of a weaved model $\wmodel$ around a Kripke model $\kmodel \in \Kwm$.
	For all $i \in [0{:}n]$, define 
	\[ Z^i = \setof{\pair {\pair {\kmodel_{\pi}} s} {\pair {\kmodel_{\pi}} s}}{ \kmodel_{\pi} \in \Kuwm \mbox{ and } |\pi| \leq (n-i)} \]
	It follows that:
	\begin{enumerate}[leftmargin=\parindent]
		\setlength{\itemsep}{0pt}
		\item $Z^{n} \subseteq Z^{(n-1)} \subseteq \dots \subseteq Z^0$ is a finite sequence of non-empty relations on $\Qof{\uwmodel} \times \Qof{\pwmodel}$ each of which satisfies the conditions (harmony), ($\R$-zig), and ($\R$-zag) in \Cref{def:wbisim};
		\item $Z^n \colon \uwmodel \bisim^n \pwmodel$; and
		\item ${\pair {\kmodel_{\pi}} s} Z^k {\pair {\kmodel_{\pi}} s}$ where $k = n - |\pi|$.
	\end{enumerate}
\end{proposition}

\lutz{it's not really clear what the motivation for this unravelling is}

\begin{corollary}\label[corollary]{prop:pwmodel}
	For all formulas $A$, if $\ideg(A) \leq n$, then, $\xforces{\uwmodel}{\pair {\kmodel_{\varepsilon}} {s_{\varepsilon}}}{A}$ iff $\xforces{\pwmodel}{\pair {\kmodel_{\varepsilon}} {s_{\varepsilon}}}{A}$.
\end{corollary}
